% 辛群（综述）
% license CCBY4
% type Wiki

本文根据 CC-BY-SA 协议转载翻译自维基百科\href{https://en.wikipedia.org/wiki/Symplectic_group}{相关文章}

在数学中，“辛群”的名称可以指两类不同但密切相关的数学群族，分别记作Sp(2n, F)与Sp(n)，其中n正整数，F为域（通常取C或 R）。后一种被称为紧致辛群，也记作$\mathrm{USp}(n)$。许多作者使用略有不同的记号，通常差别体现在系数 2 的处理上。此处采用的记号与表示这些群的最常见矩阵的规模保持一致。在卡当（Cartan）对单李代数的分类中，复群Sp(2n, C) 的李代数记作Cₙ，而Sp(n)则是Sp(2n, C)的紧致实型。需要注意的是，当我们提及“（紧致）辛群”时，通常是指依照维数n编号的一族（紧致）辛群。

“辛群”这一名称由赫尔曼·外尔提出，用来取代此前令人困惑的名称——“（线）复群”（(line) complex group）和“阿贝尔线性群”。它是“复数”一词的希腊语类比。

梅塔辛群是实数域上辛群的双覆盖；它在其他局部域、有限域以及阿德尔环上也有类似的对应物。

